\chapter{Dataflow Analysis}
\label{ch:ch09-dataflow-analysis}

An optimizer cannot trust what it sees at one program point alone.
A variable that holds zero on one path may hold a table address on another path.
A store that looks redundant in isolation may be the only store that reaches a load on some branch.
A transformation that is safe on straight-line code can therefore be unsafe once control flow joins and loops repeat.
The problem of this chapter is to compute facts that hold across control flow, conservatively and automatically, before any transformation uses them.
We solve it with techniques that propagate facts along control-flow edges until no fact changes anymore.
These techniques share one engine, one correctness argument, and one implementation shape.
They differ only in what facts they track, which direction facts flow, and how facts combine at joins.

\section{Facts that depend on paths}
\label{sec:ch09-path-facts}

A \term{dataflow analysis} computes, for every program point, a set of facts that approximates what may happen along some or all execution paths to that point.
The approximation must be conservative in a direction that keeps clients safe.
A \term{may analysis} includes every fact that can hold on an execution path, and may include additional facts because of approximation.
A \term{must analysis} reports a fact only when it holds on every path, and clients use it to act whenever the fact is guaranteed.
The choice between may and must belongs to the client, not to the solver.
Liveness is a may analysis because a variable is live when some path needs it.
Available expressions is a must analysis because an expression is available only when every path has computed it without a later kill.
Reaching definitions is a may analysis because a definition reaches when some path carries it forward.
Section~\ref{sec:ch09-equations} turns these intuitions into equations, and Section~\ref{sec:ch09-worked} computes them by hand on a small graph.
Appel develops the same client-driven view of safety for flow analyses, see~\cite{appel-1998-modern}.

Consider a function that reads a configuration flag and then computes a price.
One branch sets \code{discount} to zero and computes \code{total} from \code{price}.
The other branch sets \code{discount} from a coupon table and computes \code{total} the same way.
After the join, an optimizer wants to know whether \code{discount} is zero, whether \code{total} is redundant, and whether an earlier store to \code{total} still reaches.
No local inspection of the join block answers these questions.
The answers live on the paths, and only a whole-graph computation recovers them.
That computation is the subject of Algorithm~\ref{alg:ch09-worklist}, whose result feeds the sets in Table~\ref{tab:ch09-reaching} and the graph in Figure~\ref{fig:ch09-diamond}.

\begin{irlisting}
; B1 reads the flag, then branches to B2 or B3
; B2 computes with zero discount, B3 with table discount
; B4 joins both paths and consumes total and discount
B1: t0 = read_flag()
B1: br t0, B2, B3
B2: discount = 0
B2: total = price - discount
B2: jmp B4
B3: discount = load coupon[i]
B3: total = price - discount
B3: jmp B4
B4: store total
B4: store discount
\end{irlisting}

The listing already shows why direction matters.
Information about \code{discount} flows forward from assignments to uses.
Information about whether \code{total} is needed flows backward from uses to assignments.
A forward problem propagates facts along successor edges, while a backward problem propagates facts along predecessor edges.
The solver does not care which direction the client chose, as long as the equations name which neighbor sets feed each block.
We therefore treat direction as a parameter of one uniform framework.
The framework also explains why the join operator differs between may and must problems.
A may join takes a union because one witness path suffices, while a must join takes an intersection because every path must agree.

\section{Gen, kill, and transfer within a block}
\label{sec:ch09-gen-kill}

A \term{transfer function} summarizes what one block does to an incoming fact set.
Most classical analyses use the gen-kill form, which keeps the summary small and easy to compose.
The set \Gen{B} holds facts that block $B$ creates and that survive to the block exit.
The set \Kill{B} holds facts that block $B$ destroys on every path through the block.
The transfer function then has the shape $\Tf{B}(X) = \Gen{B} \cup (X - \Kill{B})$.
This shape says that incoming facts survive unless the block kills them, and locally created facts are added.
The form is simple enough to compute by scanning instructions once, yet expressive enough for reaching definitions, liveness, and available expressions.
Morgan presents transfer summaries in the same compositional style for optimization passes, see~\cite{morgan-1998-optimizing}.

Reaching definitions gives the cleanest example of gen and kill.
Each definition generates a fact named by its location, for example the pair of variable \code{discount} and block $B2$.
It kills every other definition of the same variable, for example the definition in $B3$.
A block that assigns \code{discount} twice generates only the last assignment, because the first dies inside the block.
A block that never assigns \code{discount} generates nothing for it and kills nothing for it.
Liveness swaps the roles of definitions and uses without changing the pattern.
Its \Gen{B} set holds variables used in $B$ before any assignment to them, and its \Kill{B} set holds variables assigned anywhere in $B$.
Available expressions uses a similar summary where a block generates an expression it computes with all operands intact, and kills an expression when it assigns to any operand.
The shared shape lets one solver serve all three clients.

A \term{basic block} contributes its summary by scanning its instructions in order.
The scan maintains a local gen set and a local kill set, both starting empty.
An assignment adds its own fact to the local gen set and removes shadowed facts from it.
It also adds every shadowed fact in the chosen finite universe to the local kill set.
A use without a preceding assignment adds the variable to the live-gen set for backward problems.
The scan visits each instruction once; its cost also includes the set operations used to construct the summaries.
The per-block summaries stay valid while the control-flow graph keeps the same instructions.
When an optimization rewrites a block, it recomputes the affected summaries and reruns the solver with valid initialization.
This locality is what makes dataflow practical inside an iterative compiler.
\section{Equations over the graph}
\label{sec:ch09-equations}

Each block satisfies two equations that relate its entry set to its exit set.
For a forward problem, the entry set combines the exits of predecessors, and the exit applies the transfer function.
For a backward problem, the exit combines the entries of successors, and the entry applies the transfer function.
We write $\In{B}$ for the entry set and $\Out{B}$ for the exit set.
We write $\Pred(B)$ for predecessors and $\Succ(B)$ for successors.
The forward equations are $\In{B} = \meet_{P \in \Pred(B)} \Out{P}$ together with $\Out{B} = \Tf{B}(\In{B})$.
The backward equations are $\Out{B} = \meet_{S \in \Succ(B)} \In{S}$ together with $\In{B} = \Tf{B}(\Out{B})$.
The operator $\meet$ is union for may analyses and intersection for must analyses.
The entry block of a forward problem and the exit blocks of a backward problem receive fixed boundary sets, usually empty, that model function entry and exit.

\begin{equation}
\Out{B} = \Gen{B} \cup (\In{B} - \Kill{B})
\label{eq:ch09-forward-transfer}
\end{equation}

Equation~\ref{eq:ch09-forward-transfer} shows the forward gen-kill transfer in isolation.
It says that exit facts are either created locally or carried through untouched.
The equation is monotone in $\In{B}$ because enlarging the input can only enlarge the output.
Monotonicity guarantees that iteration from a safe starting point cannot oscillate forever on finite sets.
Chapter 10 explains the lattice theory behind that guarantee, but the operational idea already works here.
We start every block at the identity of the join, namely empty for union and universal for intersection.
We then repeatedly apply the equations until no block changes.
The process always stops for bit-vector problems because each may set only grows and each must set only shrinks, both bounded by the finite universe of facts.
We call the limit the \term{fixed point}, and we use it as the analysis answer.
The slogan \enquote{start safe, then move monotonically} summarizes the whole method.

Joins deserve care when a block has no predecessors or no successors.
The entry block has no predecessors, so its meet ranges over an empty family.
For a may problem the empty union is the empty set, which correctly says that nothing reaches function entry from inside.
For a must problem the empty intersection is the universal set, but function entry is governed by a separate boundary equation, often with the empty set as its value.
These are different rules: an algebraic identity does not specify the program's initial state.
Prune blocks unreachable from entry before solving, so their generated facts cannot contaminate reachable joins.
For backward liveness, retain reachable loops even if they cannot reach a function exit; uses in such loops still matter.
We return to precision around unreachable code in the chapter on dead code and control dependence.

\section{One worklist solver for many analyses}
\label{sec:ch09-solver}

Algorithm~\ref{alg:ch09-worklist} solves every gen-kill problem with one worklist loop.
It initializes all block sets to the join identity, places every block on the worklist, and then processes blocks until the list empties.
Processing a block recomputes its meet from neighbors and then its transfer.
When the result differs from the stored value, the algorithm updates the stored value and enqueues the neighbors that read it.
Forward problems enqueue successors on change, while backward problems enqueue predecessors on change.
The algorithm terminates because each update moves one set strictly closer to the fixed point along a finite chain.
Upon termination every block satisfies its equations; starting at the specified extreme selects the intended extremal solution.
Soundness additionally requires conservative transfer functions and boundary facts: solving an unsound system exactly does not make it safe.
We use the word \term{worklist} for the queue itself and for the discipline of reprocessing only readers of changed facts.

\begin{algorithm}
\caption{Worklist iterative dataflow for gen-kill problems}\label{alg:ch09-worklist}
\begin{pseudocode}
\Require entry-reachable CFG $G$, gen-kill summaries, direction $d$, join $\meet$, boundary facts
\Ensure $\In{B}$ and $\Out{B}$ satisfy the dataflow equations for every $B$
\State initialize $\In{B}$ and $\Out{B}$ to the identity of $\meet$ for all $B$
\State set the boundary block to its fixed boundary value
\State $W \gets$ all blocks in $G$
\While{$W \neq \emptyset$}
\State remove some $B$ from $W$
\If{$d$ is forward}
\State $newIn \gets \meet_{P \in \Pred(B)} \Out{P}$, or boundary if $B$ is entry
\State $newOut \gets \Gen{B} \cup (newIn - \Kill{B})$
\If{$newOut \neq \Out{B}$ or $newIn \neq \In{B}$}
\State $\In{B} \gets newIn$
\State $\Out{B} \gets newOut$
\State add $\Succ(B)$ to $W$
\EndIf
\Else
\State $newOut \gets \meet_{S \in \Succ(B)} \In{S}$, or boundary if $B$ is exit
\State $newIn \gets \Gen{B} \cup (newOut - \Kill{B})$
\If{$newIn \neq \In{B}$ or $newOut \neq \Out{B}$}
\State $\In{B} \gets newIn$
\State $\Out{B} \gets newOut$
\State add $\Pred(B)$ to $W$
\EndIf
\EndIf
\EndWhile
\end{pseudocode}
\end{algorithm}

Reverse postorder processes forward problems quickly because it visits predecessors before successors along acyclic paths.
Postorder does the same service for backward problems by visiting successors first.
The cost depends on both the number of value changes and the work done on each visit, as Section~\ref{sec:ch09-implementation} explains.
An implementation should therefore invest in ordering, and use bit vectors where union, intersection, and difference are word-level operations.

A common bug is to seed must analyses with the empty set and thereby miss facts carried around loops.
Local generation can still produce nonempty answers; the error is selecting the wrong fixed point, not necessarily obtaining empty sets everywhere.
The correct seed for a must problem is the universal set of facts in the function, refined by iteration downward.
A related bug updates $\In{B}$ and $\Out{B}$ in place while still iterating over neighbors, which can hide a change from the comparison.
Algorithm~\ref{alg:ch09-worklist} avoids that trap by computing \code{newIn} and \code{newOut} into temporaries before comparing.
\section{Worked example on a diamond}
\label{sec:ch09-worked}

We now compute reaching definitions on the diamond from Figure~\ref{fig:ch09-diamond} by hand.
The universe has four facts, named $d1$ through $d4$ for the two definitions in each branch.
Block $B1$ is empty for these variables and has empty gen and kill sets.
Block $B2$ defines \code{discount} at $d1$ and \code{total} at $d2$, so its gen set is $\{d1, d2\}$ and its kill set is $\{d3, d4\}$.
Block $B3$ defines \code{discount} at $d3$ and \code{total} at $d4$, so its gen set is $\{d3, d4\}$ and its kill set is $\{d1, d2\}$.
Block $B4$ uses both variables without defining them, so its gen and kill sets are empty.
The boundary is $\In{B1} = \emptyset$, and the join is union because reaching is a may problem.
Figure~\ref{fig:visual-reaching-union} shows the facts crossing the merge, rather than only the control-flow edges.
Both definitions of a variable may reach $B4$ even though one execution took only one arm.

\begin{figure}[htbp]
  \centering
  \begin{tikzpicture}[node distance=13mm and 15mm]
    \node[basic block,align=center] (left) {$B2$ exit\\$\{d_1,d_2\}$};
    \node[basic block,align=center,right=of left] (right) {$B3$ exit\\$\{d_3,d_4\}$};
    \node[basic block,align=center,below=18mm of left,xshift=24mm] (join)
      {$B4$ entry\\$\{d_1,d_2,d_3,d_4\}$};
    \draw[cfg edge] (left) -- node[left,font=\small] {union} (join);
    \draw[cfg edge] (right) -- (join);
    \node[basic block,align=center,below=of join] (exit)
      {$B4$ exit\\$\{d_1,d_2,d_3,d_4\}$};
    \draw[cfg edge] (join) -- node[right,font=\small] {no definitions killed or generated} (exit);
  \end{tikzpicture}
  \caption{A may-analysis retains the definitions from either incoming path; the use-only join block passes them through unchanged.}
  \label{fig:visual-reaching-union}
\end{figure}


Table~\ref{tab:ch09-reaching} records the iteration, and the paragraphs below explain each pass.

\begin{figure}
\begin{tikzpicture}
\node[basic block] (b1) {B1: t0 = read flag};
\node[basic block, below left=1.2cm and 0.8cm of b1] (b2) {B2: discount = 0};
\node[basic block, below right=1.2cm and 0.8cm of b1] (b3) {B3: discount = load};
\node[basic block, below=2.6cm of b1] (b4) {B4: use total, discount};
\draw[cfg edge] (b1) -- (b2);
\draw[cfg edge] (b1) -- (b3);
\draw[cfg edge] (b2) -- (b4);
\draw[cfg edge] (b3) -- (b4);
\end{tikzpicture}
\caption{The diamond control-flow graph used for hand computation, where $B1$ branches to $B2$ and $B3$ and both join at $B4$.}
\label{fig:ch09-diamond}
\end{figure}

\begin{table}
\caption{Reaching definitions on the diamond, shown as entry and exit sets after each pass over the blocks.}
\label{tab:ch09-reaching}
\begin{tabular}{lcc}
\toprule
Block & Pass 1 entry and exit & Pass 2 entry and exit \\
\midrule
B1 & In empty, Out empty & In empty, Out empty \\
B2 & In empty, Out d1 d2 & In empty, Out d1 d2 \\
B3 & In empty, Out d3 d4 & In empty, Out d3 d4 \\
B4 & In d1 d2 d3 d4, Out same & In d1 d2 d3 d4, Out same \\
\bottomrule
\end{tabular}
\end{table}

Initialization sets every entry and exit set to empty except the boundary.
The first pass visits $B1$, computes its entry from the boundary, and keeps its exit empty.
It visits $B2$ next, reads $\In{B2} = \Out{B1} = \emptyset$, and produces $\Out{B2} = \{d1, d2\}$.
It visits $B3$ next, reads $\In{B3} = \Out{B1} = \emptyset$, and produces $\Out{B3} = \{d3, d4\}$.
It visits $B4$ last, forms $\In{B4} = \Out{B2} \cup \Out{B3} = \{d1, d2, d3, d4\}$, and copies it to $\Out{B4}$.
The second pass recomputes every block and finds no change, so the worklist empties.
The fixed point says that both definitions of \code{discount} and both definitions of \code{total} reach the join.
These facts do not justify replacing \code{discount} by zero at $B4$; other analyses might still justify restructuring the branches.
The same graph run backward for liveness would instead show that \code{price} is live on entry to both branches, a fact worth checking in the exercises.

The example also shows why union is the right join for reaching.
An intersection at $B4$ would report the empty set and claim that no definition reaches, which holds on no path.
May facts grow from empty by union, while must facts shrink from universal by intersection.

\section{Implementation notes and cost}
\label{sec:ch09-implementation}

A production implementation represents each classical fact set as a bit vector over a numbered universe.
Reaching definitions numbers every definition point, liveness numbers every variable, and available expressions numbers every expression.
Union becomes bitwise OR, intersection becomes bitwise AND, and difference becomes AND-NOT.
The per-block transfer becomes two vector operations after the one-time gen-kill scan.
A teaching compiler can stay dense throughout, while an industrial compiler tunes the representation per function size.

Let $N$ be the block count, $E$ the edge count, $U$ the fact count, and $K=\lceil U/w\rceil$ the words per bit vector on a machine with word size $w$.
The subset lattice has height $U+1$, counting elements; each stored set changes strictly at most $U$ times.
With duplicate-free queue membership, there are $O(N+EU)$ removals from the worklist, including visits that change nothing.
If $\Delta$ is the maximum number of neighbors read by a visit, recomputing every meet gives the conservative bound $O((N+EU)(\Delta+1)K)$ time after summary construction.
This bound deliberately counts repeated edge scans; a solver that propagates only newly changed bits needs a separate analysis.
Reverse postorder often reduces reprocessing by carrying facts through acyclic regions in one sweep.
Incremental solving requires an invalidation policy: deleting a definition can require removing an old may fact, which a grow-only restart from the old solution cannot do.
Reinitializing the analysis is the simple correct baseline after an edit.
Safonov discusses engineering trade-offs of this form when trusted optimization must remain predictable, see~\cite{safonov-2010-trustworthy}.

A bit-vector analysis tracks only the chosen universe, so it cannot discover that two different expressions always hold equal values.
A path-sensitive variant would separate the two branches past the join and recover that \code{discount} is zero along the left path alone.
Each step up the precision ladder costs compile time and implementation effort.
The right choice depends on which client needs the fact and how often that client runs.
The C standard leaves many evaluation and representation choices to implementations, so a portable optimizer stays conservative where the language is silent, see~\cite{wg14-2011-c11}.
% TODO(cite): need a primary reference for Kam-Ullman monotone dataflow frameworks beyond the four listed books.

\section{Exercises}
\label{sec:ch09-exercises}

\begin{enumerate}
\item \label{ex:ch09-live-diamond} Compute live variables backward on Figure~\ref{fig:ch09-diamond} by hand.
Assume that $B4$ uses \code{total}, \code{discount}, and \code{price}, that $B2$ and $B3$ define \code{total} and \code{discount}, and that neither branch uses \code{price}.
Show the gen, kill, entry, and exit sets for every block at the fixed point.
Explain in one paragraph why \code{price} is live on exit from $B1$ even though no block before $B4$ mentions it.
\item \label{ex:ch09-available-kill} Compute available expressions forward on the same diamond.
Let expression $e$ be \code{price - discount}, computed in both $B2$ and $B3$ and used in $B4$.
Show the gen and kill sets when $B2$ assigns zero to \code{discount} and $B3$ loads \code{discount} from memory.
Decide whether $e$ is available on entry to $B4$ and justify the answer from the must join.
Then change $B3$ to copy the constant zero as well and recompute the answer.
\item \label{ex:ch09-order-effect} Implement Algorithm~\ref{alg:ch09-worklist} for bit-vector reaching definitions in the language of your choice.
Run it on a chain of $N$ blocks followed by a loop back to the middle, with $N$ set to $10$, $100$, and $1000$.
Compare FIFO order, LIFO order, and reverse-postorder seeding by counting block visits to convergence.
Report the three counts for each $N$ and explain why ordering matters more as $N$ grows.
\end{enumerate}

\section{Further reading}
\label{sec:ch09-reading}

Kildall gives a foundational unified treatment of global optimization by dataflow equations; read it for the relationship between the equation system and its solution~\cite{kildall-1973-dataflow}.
Appel covers flow analysis and its use in optimization passes with implementation detail, see~\cite{appel-1998-modern}.
Morgan organizes classical optimizations around dataflow facts and transfer summaries, see~\cite{morgan-1998-optimizing}.
Safonov explains how trusted compilers bound analysis cost while keeping the safety argument explicit, see~\cite{safonov-2010-trustworthy}.
The C language definition constrains what facts an optimizer may assume across sequence points and aliasing, see~\cite{wg14-2011-c11}.
